\documentclass[10pt]{article}
\usepackage[T1]{fontenc}
\usepackage[a4paper,margin=19mm]{geometry}
\usepackage{amsmath,amssymb,booktabs,array}
\usepackage{enumitem}
\usepackage{xspace}
\usepackage{xcolor}
\usepackage[colorlinks=true,linkcolor=blue!45!black,urlcolor=blue!45!black]{hyperref}
\definecolor{accent}{HTML}{315B7D}
\setlength{\parindent}{0pt}
\setlength{\parskip}{4pt}
\setlist{nosep,leftmargin=*}
\newcommand{\goal}[1]{\medskip\noindent\textcolor{accent}{\textbf{Experimental goal.}} #1\medskip}
\newcommand{\code}[1]{\texttt{#1}}
\newcommand{\Wten}{\ifmmode\mathrm{W10}\else W10\fi}
\pagestyle{plain}

\title{TimeTensors Experiment 4: Normalization Layers}
\author{}
\date{}
\begin{document}\maketitle
\section{Theoretical overview and goal}
Normalization can improve conditioning and remove distribution shifts.  Global
standardization uses training-wide $(\mu,\sigma)$; global min--max uses training
extrema.  Instance methods instead estimate statistics from each look-back and
invert the transform after forecasting.  RevIN adds a learnable affine map to
this reversible instance transform.

\goal{Compare global and per-window alignment, distinguish standardization from
range normalization, and test whether RevIN's affine parameters add value for
DLinear and PatchTST.}

\section{Implementation details}
The methods are identity, global standard, global min--max, instance min--max,
non-affine instance normalization, and affine RevIN.  Global statistics are
computed only from accessible training look-backs to prevent leakage.  Instance
statistics are computed independently for each sample and reused to restore
predictions to data space.

The default benchmark uses six datasets, three settings, seeds 1--3, and a
DLinear table. Training defaults are random sampling, batch size 256, learning
rate $10^{-5}$, 10,000 optimizer steps, and 1,000-step validation. PatchTST is
the confirmatory second stage. The launcher is
\code{03\_normalizations.slurm}.

Report test MSE after inverse normalization, with seed standard deviation and
explicit row multiplier.  Keep the configured training loss fixed across the
normalization sweep.  Save fitted global statistics with the run so the
comparison is reproducible.

\section{Interpretation}
Global methods test optimization conditioning but cannot adapt to shifted test
levels.  Instance methods trade absolute level information for invariance.
Differences between non-affine instance normalization and RevIN isolate the
learned affine component, not the reversible transform itself.
\end{document}
